\subsection{Dyadic regularization by a genuine directed net}

\begin{theorem}[Dyadic finite Fubini--Study nets without countability]
\label{thm:fc-finite-fs-dyadic-net}
\lean{EnvelopeProperty.FSSystem.IsPSH.exists_isFinFS_net_dyadic}
\leanok
For every global psh function of rational class $\theta$ on a compact
Fubini--Study system, there is a decreasing directed net of finite rational
Fubini--Study functions with classes $\theta+2^{-n(\alpha)}\mathbf1$, where
$n$ is increasing and tends to infinity. Its pointwise infimum is the original
function. No second countability assumption is needed.
\end{theorem}
\begin{proof}
Start with the genuine decreasing finite Fubini--Study net $f_i$ whose rational
classes $a_i$ converge to $\theta$. Index the new net by all pairs $(n,i)$ with
$a_i\le\theta+2^{-n}\mathbf1$, ordered coordinatewise. For each fixed $n$, class
convergence and the strict positivity of $2^{-n}$ make this inequality hold on
an entire tail of the original directed set. Consequently the new index set
is nonempty and directed: above two admissible pairs first increase $n$ to
their maximum, then choose an old index above both old indices and above the
tail threshold for this new $n$.

Raise the class of the unchanged function $f_i$ to
$\theta+2^{-n}\mathbf1$ using Fubini--Study monotonicity. The new functions are
still decreasing because their old indices increase. The first projection
is increasing and cofinal in the natural numbers by the same tail construction,
so it tends to infinity. The second projection is cofinal in the original
index set: fix $n=0$ and choose an old index above any prescribed index and
the tail threshold. Thus the new pointwise infimum equals the unchanged
original pointwise infimum. This constructs all index data and every
approximant from the original psh function.
\end{proof}
